L'espace est muni d'un repère orthonormé \Oij.

On considère les points A$(1~;2~;~-1)$, B$(0~;3~;~2)$, C$(-2~;4~;~0)$ et D$(8~;2~;~-11)$.

\begin{enumerate}
\item
\begin{enumerate}[label=\textbf{\alph*.}]
\item Démontrer que les points A, B et C définissent un plan.
\item Démontrer que le vecteur $\vect{n}\begin{pmatrix}5\\8\\-1\end{pmatrix}$ est un vecteur normal au plan (ABC).
\item En déduire qu'une équation cartésienne du plan (ABC) est $5x+ 8y - z - 22 = 0$.
\item Vérifier que le point D n'appartient pas au plan $(ABC)$.
\end{enumerate}

\item
	\begin{enumerate}[label=\textbf{\alph*.}]
		\item On considère la droite $\Delta$, orthogonale au plan (ABC) et passant par D.

Justifier qu'une représentation paramétrique de la droite $\Delta$ est~:
\[
\begin{cases}
x = \phantom{-}8+ 5t\\
y = \phantom{-}2+ 8t\\
z = -11- t
\end{cases}, \quad \text{avec } t \in \R.
\]
		\item On appelle E le point d'intersection de la droite $\Delta$ et du plan (ABC).

Démontrer que le point E a pour coordonnées $\left(\dfrac{11}{2}~;~-2~;~-\dfrac{21}{2}\right)$.
		\item En déduire la valeur exacte de la longueur DE, distance du point D au plan (ABC).
	\end{enumerate}
\item On note H le point de coordonnées $\left(-\dfrac{2}{3}~;\dfrac{10}{3}~;\dfrac{4}{3}\right)$.

On admet qu'une représentation paramétrique de la droite (BC) est~:

\[
\begin{cases}
x = \phantom{3}-2k\\
y = 3+k\\
z = 2-2k
\end{cases}, \quad \text{avec } k \in \R.
\]

Avec un logiciel de calcul formel, on obtient les résultats suivants~:

\begin{center}
\begin{tabular}{|p{8.5cm}|}\hline
\textbf{Entrée 1~:} $(-2/3-1)*(-2)+(10/3-2)*1+(4/3-(-1))*(-2)$\\
\textbf{Sortie 1~:} $\rightarrow 0$\\
\hline
\textbf{Entrée 2~:} RésoudreSystème$(-2k=-2/3,\ 3+k=10/3,\ 2-2k=4/3)$\\
\textbf{Sortie 2~:} $\rightarrow k = 1/3$\\
\hline
\textbf{Entrée 3~:} RacineCarrée$\big((-2/3-1)^2+(10/3-2)^2+(4/3+1)^2\big)$\\
\textbf{Sortie 3~:} $\rightarrow$ RacineCarrée$(10)$\\
\hline
\textbf{Entrée 4~:} RacineCarrée$\big((-2-0)^2+(4-3)^2+(0-2)^2\big)$\\
\textbf{Sortie 4~:} $\rightarrow 3$\\
\hline
\end{tabular}
\end{center}

	\begin{enumerate}[label=\textbf{\alph*.}]
		\item En utilisant les deux premiers résultats obtenus avec la feuille de calculs ci-dessus, justifier que le point H est le projeté orthogonal du point A sur la droite (BC).
		\item En exploitant les autres résultats obtenus avec la feuille de calculs ci-dessus, déterminer le volume du tétraèdre ABCD.

On rappelle que le volume $\mathcal{V}$ d'un tétraèdre est donné par $\mathcal{V} = \dfrac{1}{3}\times\mathcal{B}\times h$, où $\mathcal{B}$ est l'aire d'une base du tétraèdre et $h$ la hauteur relative à cette base.
	\end{enumerate}
\end{enumerate}
